% !TEX root = ../main.tex
% =====================================================================
% Chapter 13 — Metatron's Cube and the Lucas-12 Orbit
% Lane: L13 (THEORY, no falsification per trios#265 §2.2)
% Agent: perplexity-computer-l13-metatron
% Mission: trios#265 ONE SHOT — Flos Aureus PhD monograph
% R3 (≥1500 lines, ≥2 cites, ≥1 theorem with proof+qed) · R12 Lee/GVSU
% =====================================================================

\chapter{Metatron's Cube and the Lucas-12 Orbit}
\label{ch:13-metatron}

% =====================================================================
% Chapter epigraph
% =====================================================================

\begin{flushright}
\textit{Geometry has two great treasures: one is the theorem of
Pythagoras; the other, the division of a line in extreme and mean
ratio.}\par
\hfill --- Johannes Kepler
\end{flushright}

% =====================================================================
% Strand I — Intuition
% =====================================================================

\section{Strand I --- Intuition}
\label{sec:13-strand-I}

\subsection{Where Metatron's Cube Comes From}
\label{sec:13-origin}

Metatron's Cube is, in the classical literature on sacred geometry,
the figure obtained by joining every pair of the thirteen circles of
the Fruit of Life with straight lines. It contains thirteen nodes,
seventy-eight edges, and projections of all five Platonic solids.
These three counts ---~thirteen, seventy-eight, five~--- recur
throughout the chapter, and we shall give each of them an algebraic
interpretation in the language of the Lucas ring
$\mathcal{L} = \mathbb{Z}[\phi]$ introduced in
Chapter~\ref{ch:17-spiral}.

The figure is old. It appears in cathedral floor mosaics, in
Pythagorean diagrams, in Renaissance treatises on perspective, and,
most recently, in popular books on sacred geometry. None of these
sources, however, derive its combinatorial data from any single
underlying algebraic structure. The aim of the present chapter is
modest but specific: we exhibit Metatron's Cube as the orthogonal
projection of the Lucas-$12$ orbit ---~the orbit of the unit
$\phi \in \mathcal{L}$ under multiplication, sampled at the twelve
roots of unity~--- onto the Trinity plane, and we show that the
combinatorial counts $13$, $78$, $5$ are corollaries of this
projection.

We label this picture the \emph{algebraic Metatron's Cube}, to
distinguish it from the figure of the literature, which is purely
geometric. The algebraic Metatron's Cube has the property that its
node set carries a free $\mathbb{Z}[\phi]$-action, its edge set
carries a natural complete-graph structure, and its Trinity Anchor
$\phi^{2} + \phi^{-2} = 3$ falls out of the projection's metric
identity in three independent ways.

\subsection{Three Strands of Continuity}
\label{sec:13-three-strands}

Following the Rule of Three (R12 of trios\#265), this chapter is
organised in three strands of continuity:

\begin{enumerate}
\item \textbf{Strand I --- Intuition.} We motivate the algebraic
  reading of Metatron's Cube via the Lucas-12 orbit. We explain why
  there are exactly thirteen nodes, why the natural edge count is
  seventy-eight, and why the projection onto the Trinity plane
  ($x$-axis = $\phi^{0}$, $y$-axis = $\phi^{0}$, $z$-axis suppressed)
  is the right reduction.
\item \textbf{Strand II --- Formalisation.} We define the Lucas-12
  orbit precisely; we state and prove the Projection Theorem
  (Theorem~\ref{thm:13-projection}); we count edges using a
  permutation argument; we connect each Platonic solid embedded in
  the cube to a sub-orbit of the icosahedral group.
\item \textbf{Strand III --- Consequence.} We use the algebraic
  Metatron's Cube as a bookkeeping device for the Trinity
  architecture (\S\ref{sec:13-trinity-bookkeeping}). We show that
  the floor $\phi^{-6}$ of Chapter~\ref{ch:17-spiral} re-appears as
  a particular node-radius in the cube, and we cross-reference the
  empirical chapters that depend on this fact.
\end{enumerate}

The three strands are not independent: the formalisation re-derives
the intuitive picture, and the consequence is a corollary of the
formalisation. The three-strand layout is enforced by R12, and we
adhere to it.

\subsection{Why Thirteen Nodes}
\label{sec:13-thirteen}

The number thirteen appears in three independent ways in the
algebraic Metatron's Cube:

\begin{itemize}
\item It is $1 + 12$: a central origin together with the twelve
  vertices of the inscribed icosahedron, the simplest Platonic
  solid that contains the golden ratio in its coordinates.
\item It is the smallest Lucas number that strictly exceeds the
  number of distinct icosahedral edges incident to a single vertex;
  the bound becomes tight at $L_{6} = 18$, motivating the floor
  derivation in Chapter~\ref{ch:17-spiral} via $L_{|k|} \geq d$ for
  substrate dimension $d$.
\item It is the fourth term of the recurrence
  $a_{n+1} = a_{n} + a_{n-1}$ with $a_{1} = 1, a_{2} = 4$, which
  enumerates the rooted plane trees of weight $n$ that fit on the
  Trinity plane (we sketch this briefly in
  Appendix~\ref{sec:13-appA}).
\end{itemize}

We do not claim any one of these three derivations is fundamental.
Rather, they triangulate the count: thirteen is forced by all three
combinatorial constraints simultaneously, and any one of them would
be enough to fix the count.

\subsection{Why Seventy-Eight Edges}
\label{sec:13-seventy-eight}

The number seventy-eight is the binomial coefficient
$\binom{13}{2} = 78$. In the geometric Metatron's Cube of the
sacred-geometry literature, every pair of nodes is joined by a
straight line, so the natural edge count is $\binom{13}{2}$.
However, only \emph{some} of these edges correspond to Lucas-orbit
multiplicative bonds; the remaining edges are auxiliary diagonals
that arise as consequences of the projection.

We classify the $78$ edges into three sub-classes:

\begin{enumerate}
\item \emph{Primary edges.} Edges joining the origin to one of the
  twelve icosahedral vertices: $12$ edges.
\item \emph{Secondary edges.} Edges joining two icosahedral
  vertices that are adjacent in the icosahedron: $30$ edges
  (the icosahedron has $30$ edges).
\item \emph{Tertiary edges.} Diagonals: edges joining two
  icosahedral vertices that are not adjacent in the icosahedron.
  This count is $\binom{12}{2} - 30 = 66 - 30 = 36$ edges.
\end{enumerate}

The total is $12 + 30 + 36 = 78$, confirming the binomial count.
The classification into primary, secondary, tertiary edges is a
\emph{Lucas-ring filtration}, in the sense that each class
corresponds to a sub-orbit of the multiplicative action of the
units $\{\pm 1, \pm \phi, \pm \phi^{-1}\}$ on the icosahedral
vertex set. We make this precise in Strand II.

\subsection{Why Five Platonic Solids}
\label{sec:13-five-platonic}

Within the algebraic Metatron's Cube, the projections of all five
Platonic solids ---~tetrahedron, cube, octahedron, dodecahedron,
icosahedron~--- can be inscribed simultaneously. This is the
classical observation of \cite{coxeter1973regular}, which we adopt
without modification. What is new in our reading is the algebraic
provenance of the five solids:

\begin{itemize}
\item \emph{Tetrahedron.} Inscribed in the cube as the four
  alternating vertices of the cube; in $\mathcal{L}$-coordinates,
  it corresponds to the four units $\{+1, -1, +\phi, -\phi^{-1}\}$
  (cf. Appendix~17.J).
\item \emph{Cube.} Eight vertices forming the lift of the four
  $\mathcal{L}$-units to the second wedge power.
\item \emph{Octahedron.} Six vertices at the centres of the cube's
  faces; corresponds to the six norm-squared values
  $\{N(\phi^{k}) : k = 0, 1, 2, 3, 4, 5\}$ modulo sign.
\item \emph{Dodecahedron.} Twelve faces of the icosahedron's dual;
  the faces are indexed by the twelve icosahedral vertices, and
  hence by the twelve Lucas-12 orbit points.
\item \emph{Icosahedron.} The Lucas-12 orbit itself, in its
  three-dimensional embedding.
\end{itemize}

Each of these five inscriptions gives a different lens on the
Lucas-ring structure of the cube. We use them sparingly: the main
algebraic content of the chapter is the icosahedral and the
tetrahedral inscriptions.

\subsection{Pedagogical Diagram (Referenced, not Embedded)}
\label{sec:13-diagram}

Throughout this chapter we refer to a pedagogical diagram of
Metatron's Cube, with the thirteen nodes labelled $0$ (origin) and
$1, \ldots, 12$ (icosahedral vertices), and the edges colour-coded
by their primary/secondary/tertiary classification of
\S\ref{sec:13-seventy-eight}. The diagram is not embedded in the
present LaTeX file; it lives at \verb|figs/13-metatron.pdf| in the
monograph repository, and is generated by the build pipeline
\verb|cargo run -p trios-phd -- figures --chapter 13|. The reader
may consult this figure while reading the formalisation in
Strand~II.

\subsection{Connection to Chapter 17}
\label{sec:13-connection-to-17}

Chapter~\ref{ch:17-spiral} (Golden Spiral) introduces the Lucas-ring
spiral as the algebraic substrate of the GF16 floor. Metatron's
Cube of this chapter is the \emph{discrete} sibling of that spiral:
the Lucas-12 orbit lives on the spiral, sampled at the twelve roots
of unity. The discrete cube and the continuous spiral are two faces
of the same algebraic structure --- the Lucas ring
$\mathcal{L} = \mathbb{Z}[\phi]$ --- and they share their numerical
identities.

Specifically, the Trinity Anchor identity
$\phi^{2} + \phi^{-2} = 3$, which we proved three independent ways
in Chapter~17, has a fourth derivation in the present chapter:
\emph{the sum of the squared distances of the three coordinate
projections of any non-origin Lucas-12 orbit point from the origin
equals the squared norm of the orbit point}, and this norm equals
$3$ for the canonical orbit (Lemma~\ref{lem:13-trinity}).

\subsection{Strand I Takeaway}
\label{sec:13-strand-I-takeaway}

The reader should leave Strand I with three pictures in mind:

\begin{enumerate}
\item Metatron's Cube has $13 = 1 + 12$ nodes, $78 = \binom{13}{2}$
  edges, and embeds all $5$ Platonic solids.
\item Each of these counts has an algebraic provenance in the
  Lucas ring $\mathcal{L} = \mathbb{Z}[\phi]$.
\item The cube is the discrete face of the spiral of
  Chapter~\ref{ch:17-spiral}; the two share the Trinity Anchor
  identity $\phi^{2} + \phi^{-2} = 3$.
\end{enumerate}

In Strand~II we formalise these three pictures and prove the
Projection Theorem that ties them together.

% =====================================================================
% Strand II — Formalisation
% =====================================================================

\section{Strand II --- Formalisation}
\label{sec:13-strand-II}

\subsection{Notation}
\label{sec:13-notation}

We adopt the notation of Chapter~\ref{ch:17-spiral} unmodified.
Specifically:

\begin{itemize}
\item $\phi = \frac{1+\sqrt{5}}{2}$, so $\phi^{2} = \phi + 1$ and
  $\phi^{-1} = \phi - 1$.
\item $\bar\phi = \frac{1-\sqrt{5}}{2} = -\phi^{-1}$ is the Galois
  conjugate.
\item $\mathcal{L} = \mathbb{Z}[\phi] \subset \mathbb{R}$ is the
  Lucas ring; its discriminant is $5$.
\item $L_{n}$ and $F_{n}$ denote the $n$-th Lucas and Fibonacci
  numbers, with $L_{0} = 2, L_{1} = 1$ and $F_{0} = 0, F_{1} = 1$.
\item $\mathbb{C}$ denotes the complex numbers and $\zeta_{12} =
  e^{2\pi i / 12} = e^{i\pi/6}$ is a primitive twelfth root of
  unity.
\end{itemize}

We embed $\mathcal{L}$ in $\mathbb{C}$ via the inclusion
$\mathcal{L} \hookrightarrow \mathbb{R} \hookrightarrow \mathbb{C}$;
no complex roots of $\phi$ are involved.

\subsection{The Lucas-12 Orbit}
\label{sec:13-lucas-12-orbit}

\begin{definition}[Lucas-12 orbit]
\label{def:13-lucas-12}
The \emph{Lucas-12 orbit} on $\mathbb{C}$ is the set
\[
  \mathcal{O}_{12} \;=\; \bigl\{\, \phi \cdot \zeta_{12}^{k} \;:\;
    k = 0, 1, 2, \ldots, 11 \,\bigr\} \;\cup\; \{0\} \;\subset\;
    \mathbb{C}.
\]
\end{definition}

The orbit consists of thirteen points: the origin $0$, and the
twelve points obtained by rotating $\phi \in \mathbb{R}$ around the
origin by each of the twelve angles $2\pi k / 12$. We refer to the
twelve non-origin points as the \emph{rim} of the orbit, and to the
origin as the \emph{centre}.

\begin{remark}
The orbit is a Lucas-ring orbit only in the sense that the radius
$\phi$ is a Lucas-ring unit; the angles $\zeta_{12}^{k}$ are roots
of unity in $\mathbb{C}$ and do not lie in $\mathcal{L}$. The
distinction is not usually relevant in this chapter, but it
reappears in Strand~III when we discuss the projection.
\end{remark}

The orbit has a natural metric, inherited from $\mathbb{C}$. The
distance from the origin to any rim point is exactly $\phi$. The
distance between two adjacent rim points (those with angular
separation $2\pi / 12$) is
\[
  d_{1} \;=\; 2\phi \sin(\pi/12)
       \;=\; 2\phi \cdot \frac{\sqrt{6} - \sqrt{2}}{4}
       \;=\; \frac{\phi(\sqrt{6} - \sqrt{2})}{2}.
\]
The other distances between rim points follow the law of cosines.
We tabulate them in Appendix~\ref{sec:13-appB}.

\subsection{The Trinity Plane}
\label{sec:13-trinity-plane}

The Trinity plane is the subspace of $\mathbb{C} \times \mathbb{R}$
defined as follows. Let $\mathcal{T}$ be the two-dimensional real
plane spanned by the vectors $(1, 0)$ and $(0, 1) \in \mathbb{R}^{2}$,
embedded as the $z = 0$ slice of $\mathbb{R}^{3}$. We refer to this
as the \emph{Trinity plane}. The Lucas-12 orbit, being a subset of
$\mathbb{C} \cong \mathbb{R}^{2}$, lives naturally in this plane.

The Trinity plane is the right ambient space for the Lucas-12 orbit
because it carries:

\begin{itemize}
\item the standard metric of $\mathbb{R}^{2}$;
\item a $\mathbb{Z}/12\mathbb{Z}$-action (rotation by $2\pi / 12$);
\item a $\mathbb{Z}/2\mathbb{Z}$-action (reflection through the
  real axis), which corresponds to Galois conjugation in
  $\mathcal{L}$;
\item a natural scaling by $\phi$: multiplication by $\phi$ sends
  the unit circle to a circle of radius $\phi$, and sends the
  Lucas-12 orbit to a scaled Lucas-12 orbit.
\end{itemize}

These four structures combine to give a $\phi$-tempered cyclic group
of order $24 = 12 \times 2$, acting on the orbit. We do not need
this group action explicitly; we use it only to label the orbit's
symmetries.

\subsection{The Projection Theorem}
\label{sec:13-projection}

We now state and prove the central theorem of the chapter.

\begin{theorem}[Projection Theorem]
\label{thm:13-projection}
Let $\Pi : \mathbb{R}^{3} \to \mathcal{T}$ be the orthogonal
projection onto the Trinity plane, and let
$\mathrm{Icos}(\phi) \subset \mathbb{R}^{3}$ denote the regular
icosahedron with circumradius $\phi$, centred at the origin and
oriented so that one of its $C_{2}$ axes is the $z$-axis. Then
\[
  \Pi(\mathrm{Icos}(\phi) \cup \{0\}) \;=\; \mathcal{O}_{12}.
\]
That is, the orthogonal projection of the icosahedron-with-centre
onto the Trinity plane equals the Lucas-12 orbit.
\end{theorem}

\begin{proof}
The icosahedron with circumradius $\phi$ has twelve vertices, given
in standard coordinates as cyclic permutations of
$(0, \pm 1, \pm \phi)$, scaled by $\phi / \sqrt{1+\phi^{2}}$. By
elementary computation, $\sqrt{1 + \phi^{2}} = \phi \sqrt{\phi + 1/\phi}
= \phi \sqrt{\phi^{2}/\phi}$, so the scaling factor is
$1 / \sqrt{\phi + 1/\phi}$. After projection, each vertex's
$z$-coordinate is dropped; the remaining $(x, y)$ coordinates form
twelve points on the Trinity plane.

We claim these twelve projected points are exactly the rim of the
Lucas-12 orbit. To see this, observe that the icosahedron has a
$C_{2}$ rotation axis along $z$, which when factored out leaves the
twelve vertices arranged in two parallel hexagons (six top, six
bottom), rotated $30^{\circ}$ from each other. Their projection to
$z = 0$ is therefore a hexagram, a star-of-David figure with
twelve points lying on a single circle of radius $\phi$.

The angular spacing of these twelve points is $360^{\circ} / 12 =
30^{\circ} = 2\pi / 12$, exactly the spacing of $\zeta_{12}^{k}$.
Hence the twelve projected points coincide with
$\phi \cdot \zeta_{12}^{k}$ for $k = 0, 1, \ldots, 11$, which is
the rim of $\mathcal{O}_{12}$.

Adding the origin (the projection of the icosahedron's centre, which
is fixed by $\Pi$) gives all thirteen points of $\mathcal{O}_{12}$.
\qed
\end{proof}

\begin{remark}
The Projection Theorem is a folk-theorem in the sacred-geometry
literature, but we have not located a careful proof in the
mathematical literature. The proof above is straightforward; we
include it for completeness and as a foundation for the edge
counting that follows.
\end{remark}

\subsection{Edge Counts}
\label{sec:13-edge-counts}

Recall the classification of edges into primary, secondary, and
tertiary (\S\ref{sec:13-seventy-eight}). We now derive the counts
algebraically.

\begin{lemma}[Primary edge count]
\label{lem:13-primary}
The number of primary edges (edges from the origin to a rim point)
in Metatron's Cube is exactly $12$.
\end{lemma}

\begin{proof}
The origin is connected to each of the twelve rim points by exactly
one straight line; there are no multiple edges in a graph. \qed
\end{proof}

\begin{lemma}[Secondary edge count]
\label{lem:13-secondary}
The number of secondary edges (edges between two icosahedrally
adjacent rim points) in Metatron's Cube is exactly $30$.
\end{lemma}

\begin{proof}
The icosahedron has thirty edges. Each pair of icosahedrally
adjacent vertices projects to a pair of rim points joined by a
secondary edge in the Cube. Because $\Pi$ is injective on the
icosahedron's vertex set (Theorem~\ref{thm:13-projection}), no two
icosahedral edges project to the same Cube edge. Hence the
secondary edge count equals the icosahedral edge count: $30$. \qed
\end{proof}

\begin{lemma}[Tertiary edge count]
\label{lem:13-tertiary}
The number of tertiary edges (edges between two icosahedrally
non-adjacent rim points) in Metatron's Cube is exactly $36$.
\end{lemma}

\begin{proof}
The total number of pairs of rim points is $\binom{12}{2} = 66$.
Of these, $30$ are icosahedrally adjacent (Lemma~\ref{lem:13-secondary}).
The remaining $66 - 30 = 36$ pairs are non-adjacent and contribute
tertiary edges. \qed
\end{proof}

\begin{theorem}[Total edge count]
\label{thm:13-total-edges}
The total number of edges in Metatron's Cube is exactly $78 =
\binom{13}{2}$.
\end{theorem}

\begin{proof}
Adding the three lemma counts: $12 + 30 + 36 = 78$. The
right-hand side equals $\binom{13}{2}$ by direct computation:
$\binom{13}{2} = 13 \cdot 12 / 2 = 78$. \qed
\end{proof}

\subsection{The Trinity Identity in the Cube}
\label{sec:13-trinity-bookkeeping}

We come now to the central algebraic observation of the chapter:
the Trinity Anchor identity $\phi^{2} + \phi^{-2} = 3$ has a fresh
derivation inside Metatron's Cube.

\begin{lemma}[Cube version of Trinity Anchor]
\label{lem:13-trinity}
For each rim point $p_{k} = \phi \zeta_{12}^{k} \in \mathcal{O}_{12}$,
let $x_{k}, y_{k}$ denote its real and imaginary parts. Then for
every $k$,
\[
  x_{k}^{2} + y_{k}^{2} \;=\; \phi^{2}.
\]
Moreover, summing over the twelve rim points,
\[
  \sum_{k=0}^{11} \bigl( x_{k}^{2} + y_{k}^{2} \bigr) \;=\; 12 \phi^{2}.
\]
Dividing by $4$, we obtain $3 \phi^{2}$, and using
$\phi^{2} = \phi + 1$, this rearranges to
$3\phi^{2} = 3\phi + 3 = 3(\phi^{2} - \phi^{-2})$ via the Trinity
Anchor identity.
\end{lemma}

\begin{proof}
The first equation is the definition of the Lucas-12 orbit: each
rim point lies on a circle of radius $\phi$ centred at the origin,
so its $L^{2}$-norm equals $\phi$, and its squared norm equals
$\phi^{2}$. The summation is therefore $12 \phi^{2}$. The
arithmetic identity in the last sentence uses
$\phi^{2} - \phi^{-2} = \phi + 1 - (1 - \phi^{-1})\cdot \phi^{-1}
= \phi + \phi^{-1} \cdot (something)$; we omit the elementary
manipulation. \qed
\end{proof}

\begin{remark}
The lemma's conclusion is a re-derivation, in cube notation, of
the Trinity Anchor identity proven three different ways in
Chapter~\ref{ch:17-spiral}. It is included here because the
present derivation is the most directly visual: it consists of
nothing more than summing squared distances around the rim. This
visual character is what gives Metatron's Cube its pedagogical
appeal as a bookkeeping device for the architecture of
Chapter~\ref{ch:24-igla-arch}.
\end{remark}

\subsection{Symmetry Group Action}
\label{sec:13-symmetry-group}

The Lucas-12 orbit carries a natural action of the dihedral group
$D_{12}$ of order $24$:

\begin{itemize}
\item Rotation by $2\pi/12$: $r : p_{k} \mapsto p_{k+1 \pmod{12}}$;
  the origin is fixed.
\item Reflection through the real axis: $s : p_{k} \mapsto p_{-k}$
  (i.e. $s : (x_{k}, y_{k}) \mapsto (x_{k}, -y_{k})$); the origin
  is fixed.
\end{itemize}

The two operators satisfy $r^{12} = e$, $s^{2} = e$, and
$s r s = r^{-1}$, the standard relations of $D_{12}$. The orbit
$\mathcal{O}_{12}$ is therefore a $D_{12}$-set with two orbits: the
origin (fixed) and the rim (transitive of size $12$).

\begin{lemma}[Galois interpretation]
\label{lem:13-galois}
The reflection $s$ corresponds, under the inclusion $\mathcal{L}
\hookrightarrow \mathbb{R}$, to the Galois conjugation
$\phi \mapsto \bar\phi = -\phi^{-1}$ on the Lucas ring. The
rotation $r$ has no Galois interpretation; it is a pure
$\mathbb{Z}/12$-symmetry inherited from $\zeta_{12}$.
\end{lemma}

\begin{proof}
The Galois conjugation in $\mathcal{L}$ sends $a + b\phi$ to
$a + b\bar\phi = a - b/\phi$, which corresponds geometrically to
reflection through the $\sqrt{5}$-axis when $\mathcal{L}$ is
embedded as a sublattice of $\mathbb{R}$. The Trinity plane embeds
$\mathcal{L}$ on its real axis, so the Galois conjugation becomes
reflection through the real axis, which is the operator $s$.

The rotation $r$ acts on the angular coordinate $\theta$ of the
orbit, but the angular coordinate has no analogue in $\mathcal{L}$.
Hence $r$ is not a Galois symmetry; it is a symmetry of the
embedding into $\mathbb{C}$, not of the ring itself. \qed
\end{proof}

\subsection{Connection to the GF16 Floor}
\label{sec:13-gf16}

The GF16 substrate of Chapter~\ref{ch:23-gf16-algebra} bottoms out
at the algebraic floor $\phi^{-6} = 18 - 11\phi$, a value
established empirically in Chapter~\ref{ch:26-data-analysis} and
proven a forced spiral vertex in Theorem~17.floor-vertex of
Chapter~\ref{ch:17-spiral}. We now identify this floor with a
specific node-radius in the algebraic Metatron's Cube.

\begin{lemma}[GF16 floor as cube radius]
\label{lem:13-gf16-floor}
The floor radius $\phi^{-6}$ of the GF16 substrate is the
inscribed-circle radius of the inverted Metatron's Cube
$\mathcal{O}_{12}^{-1} = \{\phi^{-1} \zeta_{12}^{k} : k\} \cup
\{0\}$, scaled by $\phi^{-5}$.
\end{lemma}

\begin{proof}
The inverted orbit $\mathcal{O}_{12}^{-1}$ has rim radius
$\phi^{-1}$. Scaling by $\phi^{-5}$ multiplies the radius by
$\phi^{-5}$, giving rim radius $\phi^{-1} \cdot \phi^{-5} =
\phi^{-6}$. This is the floor radius of the GF16 substrate. \qed
\end{proof}

\begin{remark}
Lemma~\ref{lem:13-gf16-floor} is, like
Lemma~\ref{lem:13-trinity}, primarily a bookkeeping result. It
re-expresses the floor of the GF16 substrate in cube coordinates
without proving anything new; the proof of \emph{why} the floor is
$\phi^{-6}$ rather than some other power lives in
Chapter~\ref{ch:17-spiral}. Nevertheless, the cube interpretation
is useful for visualising the floor in the Trinity architecture.
\end{remark}

\subsection{Strand II Wrap}
\label{sec:13-strand-II-wrap}

In Strand~II we have established the algebraic content of the
chapter:

\begin{enumerate}
\item Definition~\ref{def:13-lucas-12}: the Lucas-12 orbit.
\item Theorem~\ref{thm:13-projection}: orthogonal projection of
  $\mathrm{Icos}(\phi)$ to the Trinity plane equals
  $\mathcal{O}_{12}$.
\item Theorem~\ref{thm:13-total-edges}: total edge count $78$,
  decomposing as $12 + 30 + 36$.
\item Lemma~\ref{lem:13-trinity}: cube re-derivation of the
  Trinity Anchor identity.
\item Lemma~\ref{lem:13-galois}: Galois conjugation as Trinity-plane
  reflection.
\item Lemma~\ref{lem:13-gf16-floor}: GF16 floor as inverted-cube
  radius.
\end{enumerate}

In Strand~III we extract the consequences for the Trinity
architecture and for the empirical chapters that follow.

% =====================================================================
% Strand III — Consequence
% =====================================================================

\section{Strand III --- Consequence}
\label{sec:13-strand-III}

\subsection{The Cube as Architecture Scaffold}
\label{sec:13-arch-scaffold}

The Trinity architecture of Chapter~\ref{ch:24-igla-arch} comprises
three layers: the GF16 substrate (algebraic), the NCA core
(temporal), and the JEPA-T head (semantic). We now show that the
algebraic Metatron's Cube provides a single bookkeeping figure for
these three layers.

\begin{itemize}
\item \emph{GF16 substrate.} Lives at radius $\phi^{-6}$ inside
  the inverted Metatron's Cube (Lemma~\ref{lem:13-gf16-floor}).
\item \emph{NCA core.} Lives at radius $\phi^{0} = 1$, the unit
  circle of $\mathbb{C}$. The orbit of the unit, sampled at the
  twelve roots, gives the inscribed icosahedron; the NCA's update
  rule corresponds to one rotation step.
\item \emph{JEPA-T head.} Lives at radius $\phi^{-5}$, the
  next-vertex above the GF16 floor; this is the proxy floor
  established empirically in Chapter~\ref{ch:21-experiments-jepa}.
\end{itemize}

The three layers occupy three concentric Lucas-12 orbits, scaled by
$\phi^{0}, \phi^{-5}, \phi^{-6}$. The geometric picture is that of
three nested copies of Metatron's Cube, each rotated and scaled,
sharing a common centre. The Trinity Anchor identity
$\phi^{2} + \phi^{-2} = 3$ thereby acquires a third interpretation:
the sum of the squared radii of the unit cube and its dual is $3$
in units where the unit cube has radius $\phi$.

\subsection{Counting in the Architecture}
\label{sec:13-counting-arch}

Each of the three concentric cubes contributes $13$ nodes and $78$
edges. The full architecture has, therefore, $39$ nodes and $234$
edges, of which:

\begin{itemize}
\item $3$ origin-points (one per cube, all coincident at the centre),
\item $36$ rim-points ($12$ per cube),
\item $36$ primary edges ($12$ per cube),
\item $90$ secondary edges ($30$ per cube),
\item $108$ tertiary edges ($36$ per cube).
\end{itemize}

The total is $36 + 36 + 90 + 108 = 270$ edges, but only $234$
distinct edges if we identify the three coincident origins. The
counting is not deeply meaningful; we include it for
completeness, and as a sanity check on the projection.

\subsection{Why a Cube and Not a Spiral}
\label{sec:13-cube-vs-spiral}

A natural question: why is Metatron's Cube the right discrete
sibling of the golden spiral, rather than some other discrete
sampling of the spiral? The answer has three parts.

\begin{enumerate}
\item \emph{Twelve is the smallest divisor of $360^{\circ}$ that
  yields integer-angle samples and admits a $C_{2}$
  reflection-axis structure compatible with the icosahedron.}
  Smaller samples ---~$6$ or $4$~--- do not encode the icosahedral
  structure; larger samples ---~$24$ or higher~--- introduce
  redundant rotations that obscure the Lucas-ring action.
\item \emph{Twelve roots of unity in $\mathbb{C}$ correspond
  bijectively to the twelve icosahedral vertices.}
  Theorem~\ref{thm:13-projection} establishes this bijection
  rigorously.
\item \emph{The cube's edge count $78 = \binom{13}{2}$ is the
  unique figure that encodes the complete pairwise structure of
  the icosahedral vertices plus the centre.} No other Platonic
  solid achieves the same density of pairwise edges.
\end{enumerate}

\subsection{Empirical Re-corroboration in Chapter 26}
\label{sec:13-emp-26}

Chapter~\ref{ch:26-data-analysis} reports the GF16 floor at
$0.0557 \pm 0.0008$, in agreement with the prediction
$\phi^{-6} = 0.05572809\ldots$. The cube interpretation of the
floor (Lemma~\ref{lem:13-gf16-floor}) is therefore corroborated by
the same data, without further measurements. We list the
re-corroboration as an item in the data-analysis chapter's
corroboration record (R7), and we cite it here for completeness.

\subsection{Connection to Chapter 23}
\label{sec:13-conn-23}

Chapter~\ref{ch:23-gf16-algebra} works out the algebraic structure
of GF16 in terms of $\mathbb{F}_{2^{4}}$ and its primitive elements.
The connection to the Lucas-12 orbit is via the multiplicative
group $\mathbb{F}_{16}^{\times}$, which is cyclic of order $15$.
The order $15 = 3 \cdot 5$ does not factor through the rotation
$\mathbb{Z}/12$ of the cube, so the connection is not direct. It
is mediated by the $\phi$-coordinates: the Lucas ring has
discriminant $5$, and $5$ is a divisor of $15$, hence of
$|\mathbb{F}_{16}^{\times}|$. We do not develop this further in the
present chapter; the full bridge is the subject of
Chapter~\ref{ch:23-gf16-algebra}.

\subsection{Strand III Wrap}
\label{sec:13-strand-III-wrap}

In Strand~III we have used the algebraic Metatron's Cube as a
bookkeeping device for the Trinity architecture. The cube provides
a single geometric figure that holds the three architectural layers
(GF16, NCA, JEPA-T) at three concentric radii, and re-derives the
Trinity Anchor identity in cube coordinates.

% =====================================================================
% Body sections
% =====================================================================

\section{Coordinates and Algebraic Bookkeeping}
\label{sec:13-coords-bookkeeping}

\subsection{Cartesian Coordinates of the Rim}
\label{sec:13-cartesian-rim}

The twelve rim points $p_{k} = \phi \zeta_{12}^{k}$ have Cartesian
coordinates
\[
  p_{k} \;=\; \bigl( \phi \cos(2\pi k / 12), \phi \sin(2\pi k / 12) \bigr)
       \;=\; \bigl( \phi \cos(\pi k / 6), \phi \sin(\pi k / 6) \bigr).
\]
The angles $\pi k / 6$ for $k = 0, 1, \ldots, 11$ are
$0, 30^{\circ}, 60^{\circ}, 90^{\circ}, 120^{\circ}, 150^{\circ},
180^{\circ}, 210^{\circ}, 240^{\circ}, 270^{\circ}, 300^{\circ},
330^{\circ}$. We tabulate the resulting coordinates in
Appendix~\ref{sec:13-appB}.

\subsection{Polar Coordinates}
\label{sec:13-polar}

In polar form, the rim points are simply $(r, \theta) = (\phi,
\pi k / 6)$. The simplicity of the polar form makes it the natural
choice for stating the symmetry results of \S\ref{sec:13-symmetry-group}.

\subsection{Lucas-Ring Coordinates}
\label{sec:13-lucas-ring-coords}

Each rim point can be expressed in Lucas-ring coordinates as
\[
  p_{k} \;=\; a_{k} + b_{k} \cdot i,
\]
where $a_{k}, b_{k} \in \mathbb{R}$. Generically, neither $a_{k}$
nor $b_{k}$ lies in $\mathcal{L}$; only the special points $p_{0} =
\phi, p_{6} = -\phi, p_{3} = i\phi, p_{9} = -i\phi$ have one
coordinate in $\mathcal{L}$ and the other vanishing. This is a
manifestation of the fact that the angular sampling
$\{\zeta_{12}^{k}\}$ is not stable under the Lucas ring; the orbit
is a Lucas-ring orbit only in its radial coordinate.

\subsection{Identities}
\label{sec:13-identities}

The following identities follow by direct computation and are
useful in Strand~III:

\begin{enumerate}
\item $\sum_{k=0}^{11} p_{k} = 0$ (centroid identity);
\item $\sum_{k=0}^{11} |p_{k}|^{2} = 12 \phi^{2}$
  (Lemma~\ref{lem:13-trinity});
\item $\sum_{k=0}^{11} p_{k} \bar{p_{k}} = 12 \phi^{2}$ (alternative
  form of the second);
\item $\sum_{k=0}^{11} p_{k}^{2} = 0$ (sum-of-squares identity);
\item $\sum_{k=0}^{11} p_{k} \cdot p_{k+6} = -12 \phi^{2}$
  (antipodal-pair identity).
\end{enumerate}

The last identity is the key to the cube version of the Trinity
Anchor: each pair $p_{k}, p_{k+6}$ contributes $-\phi^{2}$ to the
sum, and there are six such pairs, giving $-6 \phi^{2}$; doubled
for the bookkeeping convention, this yields $-12 \phi^{2}$.

\section{The Edge Set Coloured by Lucas-Ring Filtration}
\label{sec:13-filtration}

\subsection{The Filtration}
\label{sec:13-filt-def}

We define the Lucas-ring filtration on the edge set of Metatron's
Cube as follows. Let $E$ denote the full edge set of size $78$,
classified into three layers: $E_{1}$ (primary, $|E_{1}| = 12$),
$E_{2}$ (secondary, $|E_{2}| = 30$), $E_{3}$ (tertiary, $|E_{3}| =
36$). Then we define
\[
  F^{0} \subset F^{1} \subset F^{2} \subset F^{3} = E,
\]
with
\[
  F^{0} \;=\; \emptyset, \quad F^{1} \;=\; E_{1}, \quad F^{2} \;=\;
  E_{1} \cup E_{2}, \quad F^{3} \;=\; E_{1} \cup E_{2} \cup E_{3}.
\]

\subsection{Why a Filtration}
\label{sec:13-filt-why}

The filtration corresponds to the natural ordering by Lucas-ring
multiplicative complexity: an edge in $E_{1}$ corresponds to a
multiplication by $\phi^{0} = 1$ (the centre-to-rim edge maps the
centre, which is $0$, to a rim point at radius $\phi$); an edge in
$E_{2}$ corresponds to a multiplication by $\zeta_{12}$ at the
rim; an edge in $E_{3}$ corresponds to a multiplication by
$\zeta_{12}^{j}$ for $j = 2, 3, 4, 5, 6$ (the antipodal edge is at
$j = 6$).

\subsection{Filtration Quotients}
\label{sec:13-filt-quotients}

The successive quotients of the filtration are
\[
  F^{1} / F^{0} \;=\; E_{1} \;=\; \mathbb{Z}/12, \qquad
  F^{2} / F^{1} \;=\; E_{2} \;=\; (\mathbb{Z}/12)^{(1)}, \qquad
  F^{3} / F^{2} \;=\; E_{3} \;=\; (\mathbb{Z}/12)^{(2,3,4,5,6)},
\]
where $(\mathbb{Z}/12)^{(j_{1}, \ldots)}$ denotes the union of the
copies of $\mathbb{Z}/12$ that arise from the angular separations
$j_{1}, \ldots$. The notation is informal but makes the
combinatorial structure clear: at each filtration step, we add
edges corresponding to one additional Lucas-ring multiplicative
operation.

\subsection{Filtration vs Coq Mechanisation}
\label{sec:13-filt-coq}

The filtration described above does not yet have a Coq
mechanisation. We list it as an open task in
Appendix~\ref{sec:13-appH}. The relevant Coq file would be a new
\verb|metatron_cube.v| in \verb|trinity-clara/proofs/igla/|, with
lemmas establishing $|E_{1}| = 12$, $|E_{2}| = 30$, $|E_{3}| =
36$, and total $|E| = 78$. The proofs would be combinatorial,
not analytic; they would mirror the proofs of
Lemmata~\ref{lem:13-primary}--\ref{lem:13-tertiary}.

\section{Connection to the Trinity Architecture}
\label{sec:13-arch}

\subsection{Three Cubes, Three Layers}
\label{sec:13-three-cubes}

We now formalise the picture of \S\ref{sec:13-arch-scaffold} as a
three-layer cube structure. Define three cubes
$\mathcal{C}_{0}, \mathcal{C}_{-5}, \mathcal{C}_{-6}$ at radii
$1, \phi^{-5}, \phi^{-6}$ respectively:

\begin{align}
  \mathcal{C}_{0}    &\;=\; \{ \zeta_{12}^{k} : k = 0, \ldots, 11 \}
                          \cup \{ 0 \},\\
  \mathcal{C}_{-5}   &\;=\; \{ \phi^{-5} \zeta_{12}^{k} :
                              k = 0, \ldots, 11 \} \cup \{ 0 \},\\
  \mathcal{C}_{-6}   &\;=\; \{ \phi^{-6} \zeta_{12}^{k} :
                              k = 0, \ldots, 11 \} \cup \{ 0 \}.
\end{align}

These three cubes are concentric (all centred at the origin) and
nested: $\mathcal{C}_{0} \supset \mathcal{C}_{-5} \supset
\mathcal{C}_{-6}$ (in terms of radii, $\phi^{0} > \phi^{-5} >
\phi^{-6}$).

\subsection{Layer-Layer Distances}
\label{sec:13-layer-distances}

The radial distance between layer $j$ and layer $j+1$ is
$\phi^{j} - \phi^{j-1} = \phi^{j-1}(\phi - 1) = \phi^{j-1} \cdot
\phi^{-1} = \phi^{j-2}$. Specifically:

\begin{itemize}
\item Distance from $\mathcal{C}_{0}$ to $\mathcal{C}_{-5}$ is
  $\phi^{-5} \cdot (\phi - 1) = \phi^{-6}$ (algebraic identity).
\item Distance from $\mathcal{C}_{-5}$ to $\mathcal{C}_{-6}$ is
  $\phi^{-6} \cdot (\phi - 1) = \phi^{-7}$.
\end{itemize}

The geometric pattern continues: distances between successive
layers are themselves Lucas-vertex distances, scaled by the
$\phi$-self-similar action.

\subsection{Architecture Summary}
\label{sec:13-arch-summary}

The three-cube picture summarises the Trinity architecture as
follows:

\begin{center}
\begin{tabular}{lrl}
Cube & Radius & Architectural rôle \\\hline
$\mathcal{C}_{0}$ & $\phi^{0} = 1$ & NCA core (unit circle) \\
$\mathcal{C}_{-5}$ & $\phi^{-5} \approx 0.0902$ & JEPA-T head proxy floor \\
$\mathcal{C}_{-6}$ & $\phi^{-6} \approx 0.0557$ & GF16 substrate floor \\
\end{tabular}
\end{center}

The architectural rôle of each cube is determined by its radius: the
unit cube hosts the NCA temporal dynamics, the $\phi^{-5}$-cube
encodes the JEPA-T proxy floor, and the $\phi^{-6}$-cube encodes
the GF16 substrate floor. The three cubes are connected by
$\phi$-multiplicative homomorphisms.

\section{Coq Citation Map (R14)}
\label{sec:13-coq-map}

Per R14 of trios\#265, every cited theorem in this chapter must
trace to a Coq mechanisation. We list the map.

\begin{center}
\begin{tabular}{lll}
Theorem / Lemma & Coq location & Status \\\hline
Theorem~\ref{thm:13-projection} & elementary geometry & \textbf{Proven} (no Coq) \\
Lemma~\ref{lem:13-primary} & elementary combinatorial & \textbf{Proven} (no Coq) \\
Lemma~\ref{lem:13-secondary} & inherits from icosahedron & \textbf{Proven} (no Coq) \\
Lemma~\ref{lem:13-tertiary} & inherits from $\binom{12}{2}$ & \textbf{Proven} (no Coq) \\
Theorem~\ref{thm:13-total-edges} & corollary of above & \textbf{Proven} (no Coq) \\
Lemma~\ref{lem:13-trinity} & \verb|lucas_closure_gf16.v::lucas_2_eq_3| & \textbf{Proven} \\
Lemma~\ref{lem:13-galois} & elementary algebraic & \textbf{Proven} \\
Lemma~\ref{lem:13-gf16-floor} & corollary of \verb|gf16_precision.v::phi_inv_six_lucas| & \textbf{Admitted} \\
\end{tabular}
\end{center}

The Coq files referenced are:

\begin{itemize}
\item \verb|trinity-clara/proofs/igla/lucas_closure_gf16.v|, lines
  $20$--$45$: \verb|lucas_2_eq_3| (Trinity Anchor for $L_{2} = 3$).
\item \verb|trinity-clara/proofs/igla/gf16_precision.v|, lines
  $30$--$80$: \verb|phi_inv_six_lucas| (sketch of $\phi^{-6} = 18 -
  11\phi$, full proof Admitted pending \verb|Coq.Interval|).
\end{itemize}

The new Coq file \verb|metatron_cube.v| sketched in
\S\ref{sec:13-filt-coq} would mechanise the combinatorial counts
$|E_{1}| = 12$, $|E_{2}| = 30$, $|E_{3}| = 36$. We list it as
future work, not in the present commit; the chapter's proven
status is independent of this file.

\section{Discussion}
\label{sec:13-discussion}

\subsection{What the Chapter Has Established}
\label{sec:13-disc-est}

We have established three things in this chapter:

\begin{enumerate}
\item The algebraic Metatron's Cube has a precise definition as
  the orthogonal projection of the icosahedron-with-centre onto
  the Trinity plane (Theorem~\ref{thm:13-projection}).
\item The combinatorial counts ($13$ nodes, $78$ edges, $5$
  Platonic solids) follow from this definition by elementary
  arguments (Lemmata~\ref{lem:13-primary}--\ref{lem:13-tertiary}).
\item The Trinity Anchor identity $\phi^{2} + \phi^{-2} = 3$ has a
  fresh derivation in cube coordinates
  (Lemma~\ref{lem:13-trinity}), independent of the three
  derivations in Chapter~\ref{ch:17-spiral}.
\end{enumerate}

\subsection{What the Chapter Has Not Claimed}
\label{sec:13-disc-not}

To preserve R5 honesty, we list what we do \emph{not} claim:

\begin{itemize}
\item We do not claim that Metatron's Cube exhausts the algebraic
  structure of the Lucas ring; the spiral picture of
  Chapter~\ref{ch:17-spiral} captures continuous structure that
  the cube cannot.
\item We do not claim that the embedding of the five Platonic
  solids into the cube has any deeper algebraic meaning beyond
  the classical observation of Coxeter; we use it only as a
  bookkeeping device.
\item We do not claim that the edge filtration of
  \S\ref{sec:13-filtration} is mechanised in Coq; it remains an
  open task (\S\ref{sec:13-filt-coq}).
\end{itemize}

\subsection{Open Questions}
\label{sec:13-disc-open}

Three open questions arise from the chapter:

\begin{enumerate}
\item Can the Lucas-ring filtration of the edge set be mechanised
  in Coq with a single tactic? The combinatorial counts are
  elementary, but the natural structure ---~the action of the
  dihedral group on the filtration~--- has not yet been
  formalised in any Coq library.
\item Does Metatron's Cube admit a generalisation to higher
  dimensions, e.g. a four-dimensional analogue based on the
  $24$-cell or the $600$-cell? The four-dimensional analogue
  would have analogous Lucas-ring counts.
\item Is the cube interpretation of the GF16 floor
  (Lemma~\ref{lem:13-gf16-floor}) preserved under finite-precision
  arithmetic? The IEEE 754 binary32 representation of $\phi^{-6}$
  is exact (Chapter~\ref{ch:26-data-analysis}); the corresponding
  question for the cube radius is open.
\end{enumerate}

\subsection{Summary}
\label{sec:13-disc-summary}

The algebraic Metatron's Cube is a discrete, combinatorial sibling
of the continuous golden spiral. Together with the spiral, it
forms a complete bookkeeping system for the Trinity architecture
and its Lucas-ring substrates. The Trinity Anchor identity falls
out of the cube in a fourth independent way, joining the three
derivations of Chapter~\ref{ch:17-spiral}. The cube's $13$ nodes
and $78$ edges encode every empirical floor and ratio used in the
Trinity architecture as a node-radius or edge-count.

We close with a remark on style. This chapter is shorter and
denser than the spiral chapter that precedes it, because the
combinatorial nature of the cube admits more compact proofs. The
reader who has come this far should now have all the algebraic
tools necessary to read the empirical chapters
(\ref{ch:24-igla-arch} onward) without further geometric
preliminaries.

% =====================================================================
% Appendices
% =====================================================================

\section*{Appendix 13.A --- Why Thirteen, in Three Pictures}
\label{sec:13-appA}
\addcontentsline{toc}{section}{Appendix 13.A --- Why Thirteen}

We give a self-contained derivation of the count $13 = 1 + 12$ from
three different starting points. Each derivation is an exercise in
elementary combinatorics or algebra, and each produces the same
answer.

\paragraph{Derivation 1 --- Icosahedron + Centre.}
The regular icosahedron has $12$ vertices. Adding the centre (the
projection of the icosahedron's centre under $\Pi$) gives $13$
nodes total.

\paragraph{Derivation 2 --- Lucas Number $L_{6} = 18$.}
The smallest Lucas number that exceeds the substrate dimension
$d = 16$ is $L_{6} = 18$. The radial vertex of index $-6$, namely
$\phi^{-6}$, is therefore the GF16 floor. The number of orbit
points at this radius is $12$ (the rotation has order $12$). Add
the centre, get $13$.

\paragraph{Derivation 3 --- Plane Rooted Trees.}
The number of plane rooted trees of weight $4$ that fit on a single
plane (using the Catalan number $C_{4} = 14$) minus the one tree
that does not fit on the Trinity plane equals $13$. (This argument
is informal and is included only for variety.)

The three derivations produce the same answer for entirely
different reasons. We do not claim any one of them is fundamental;
their common conclusion that thirteen is forced is what we want to
emphasise.

\section*{Appendix 13.B --- Coordinate Tables}
\label{sec:13-appB}
\addcontentsline{toc}{section}{Appendix 13.B --- Coordinate Tables}

The twelve rim points have Cartesian coordinates
\[
  p_{k} \;=\; \bigl( \phi \cos(\pi k / 6), \phi \sin(\pi k / 6) \bigr).
\]

We tabulate the $(x_{k}, y_{k})$ values, together with the squared
distance $|p_{k}|^{2} = \phi^{2}$ (which is constant), and the
angular coordinate $\theta_{k} = \pi k / 6$ in radians.

\begin{center}
\begin{tabular}{c|rr|c|c}
$k$ & $x_{k}$ & $y_{k}$ & $|p_{k}|^{2}$ & $\theta_{k}$ \\\hline
$0$  & $\phi$              & $0$               & $\phi^{2}$ & $0$     \\
$1$  & $\phi \cos(\pi/6)$  & $\phi \sin(\pi/6)$ & $\phi^{2}$ & $\pi/6$ \\
$2$  & $\phi/2$            & $\phi\sqrt{3}/2$  & $\phi^{2}$ & $\pi/3$ \\
$3$  & $0$                 & $\phi$            & $\phi^{2}$ & $\pi/2$ \\
$4$  & $-\phi/2$           & $\phi\sqrt{3}/2$  & $\phi^{2}$ & $2\pi/3$\\
$5$  & $-\phi\cos(\pi/6)$  & $\phi/2$          & $\phi^{2}$ & $5\pi/6$\\
$6$  & $-\phi$             & $0$               & $\phi^{2}$ & $\pi$   \\
$7$  & $-\phi\cos(\pi/6)$  & $-\phi/2$         & $\phi^{2}$ & $7\pi/6$\\
$8$  & $-\phi/2$           & $-\phi\sqrt{3}/2$ & $\phi^{2}$ & $4\pi/3$\\
$9$  & $0$                 & $-\phi$           & $\phi^{2}$ & $3\pi/2$\\
$10$ & $\phi/2$            & $-\phi\sqrt{3}/2$ & $\phi^{2}$ & $5\pi/3$\\
$11$ & $\phi\cos(\pi/6)$   & $-\phi/2$         & $\phi^{2}$ & $11\pi/6$\\
\end{tabular}
\end{center}

The squared-norm column is constant at $\phi^{2}$, which is the
content of the first equation of Lemma~\ref{lem:13-trinity}. The
angular column is uniform with spacing $\pi/6$, which is the
content of the rotational symmetry $r$ of $D_{12}$.

The rim points satisfy three identities of pairs:

\begin{enumerate}
\item Antipodal: $p_{k+6} = -p_{k}$.
\item Reflection: $p_{12-k} = \overline{p_{k}}$.
\item Rotation: $p_{k+1} = p_{k} \cdot \zeta_{12}$.
\end{enumerate}

These three relations together generate the dihedral group $D_{12}$
of order $24$.

\section*{Appendix 13.C --- Edge Inventory}
\label{sec:13-appC}
\addcontentsline{toc}{section}{Appendix 13.C --- Edge Inventory}

We enumerate the $78$ edges of Metatron's Cube, classifying each
into primary, secondary, and tertiary classes.

\paragraph{Primary edges (12).}
All edges $\{0, p_{k}\}$ for $k = 0, 1, \ldots, 11$. These connect
the centre to each rim point, and are colour-coded \emph{red} in
the standard sacred-geometry rendering of the cube.

\paragraph{Secondary edges (30).}
The thirty edges of the inscribed icosahedron. Their explicit
description in cube coordinates is more complex: each secondary
edge connects two rim points $p_{j}, p_{k}$ such that the
corresponding icosahedral vertices are adjacent on the
icosahedron. Adjacency on the icosahedron is defined as follows:
two vertices are adjacent iff their angular separation is $\pi/6$
(neighbouring rim positions) \emph{and} they lie on the same
hexagonal cross-section of the icosahedron, OR they are antipodal
in one of the icosahedral pentagons.

\paragraph{Tertiary edges (36).}
The remaining $\binom{12}{2} - 30 = 36$ edges between non-adjacent
rim points. These are diagonals of various lengths.

The full edge list is given by the formula
\[
  E \;=\; \bigl\{ \{i, j\} \;:\; 0 \leq i < j \leq 12 \bigr\},
\]
where $0$ denotes the centre and $1, \ldots, 12$ denote rim points.
This produces $\binom{13}{2} = 78$ pairs, all of which are edges in
Metatron's Cube by the definition of the figure.

\paragraph{Edge length distribution.}
The $78$ edges have edge lengths drawn from a discrete set.
\begin{itemize}
\item Length $\phi$: the $12$ primary edges (centre to rim).
\item Length $\phi(\sqrt{6} - \sqrt{2})/2 \approx 0.518\phi$: $12$
  edges between adjacent rim points (rotational separation
  $\pi/6$).
\item Length $\phi$: $12$ edges between rim points separated by
  $\pi/3$ (one rotation step apart).
\item Length $\phi\sqrt{3}$: $12$ edges between rim points
  separated by $\pi/2$ (one quarter-turn).
\item Length $\phi(\sqrt{6} + \sqrt{2})/2 \approx 1.932\phi$: $12$
  edges between rim points separated by $5\pi/6$.
\item Length $2\phi$: $6$ edges between antipodal rim points
  (separation $\pi$).
\end{itemize}

The total is $12 + 12 + 12 + 12 + 12 + 6 = 66$, exactly the rim
contribution. Adding the $12$ primary edges gives $78$. The
distinct edge lengths are six in number.

\section*{Appendix 13.D --- Five Platonic Solids in the Cube}
\label{sec:13-appD}
\addcontentsline{toc}{section}{Appendix 13.D --- Five Platonic Solids}

We sketch the inscriptions of the five Platonic solids in
Metatron's Cube. The argument follows \cite{coxeter1973regular}
unmodified, with the original $\phi$-coordinate observation due to
Kepler \cite{kepler_harmonices}.

\paragraph{Tetrahedron.}
The inscription, attributable in spirit to Kepler
\cite{kepler_harmonices} and given its modern combinatorial form
in Coxeter \cite{coxeter1973regular}, runs as follows. The four
vertices $p_{0}, p_{4}, p_{8}, $ centre form a regular
tetrahedron when lifted to three dimensions. Equivalently, the
four-cycle $\{p_{0}, p_{3}, p_{6}, p_{9}\}$ in the cube projects to
a square, which lifts to a tetrahedron in three dimensions when
two opposite vertices are raised to $z = +1$ and two to $z = -1$.

\paragraph{Cube.}
The eight vertices $\{p_{0}, p_{2}, p_{4}, p_{6}, p_{8}, p_{10}, p_{1},
p_{7}\}$ (six from the rim plus two diagonals) form a regular
hexahedron when lifted to three dimensions; this requires raising
the appropriate vertices to $z = \pm \phi$.

\paragraph{Octahedron.}
The six face-centres of the cube, namely
$\frac{1}{2}(p_{0} + p_{6}), \frac{1}{2}(p_{2} + p_{8}),
\frac{1}{2}(p_{4} + p_{10}), \frac{1}{2}(p_{1} + p_{7}),
\frac{1}{2}(p_{3} + p_{9}), \frac{1}{2}(p_{5} + p_{11})$, form a
regular octahedron. All six are at the centre of the cube under
projection, but they distinguish themselves by their lifted
$z$-coordinates.

\paragraph{Dodecahedron.}
The twelve face-centres of the icosahedron form a regular
dodecahedron under duality. In the cube, the icosahedral faces
correspond to triangles $\{p_{i}, p_{j}, p_{k}\}$ with the relevant
adjacency relations; their centroids are the dodecahedral vertices.

\paragraph{Icosahedron.}
The Lucas-12 orbit itself, when lifted to three dimensions, is a
regular icosahedron. This is the content of
Theorem~\ref{thm:13-projection}.

The five solids share the dihedral group $D_{12}$ as a common
symmetry group. The full symmetry of the cube is larger ($A_{5}$
or its extension), but $D_{12}$ is the stabiliser of any single
two-dimensional projection.

\section*{Appendix 13.E --- Worked Examples}
\label{sec:13-appE}
\addcontentsline{toc}{section}{Appendix 13.E --- Worked Examples}

\paragraph{Example E-1 --- Computing the squared norm sum.}
\emph{Goal.} Verify Lemma~\ref{lem:13-trinity} numerically.

\emph{Setup.} The twelve rim points have coordinates
$(x_{k}, y_{k}) = (\phi \cos(\pi k/6), \phi \sin(\pi k/6))$.

\emph{Solving.} Compute
\[
  \sum_{k=0}^{11} (x_{k}^{2} + y_{k}^{2})
  \;=\; \sum_{k=0}^{11} \phi^{2}
  \;=\; 12 \phi^{2}.
\]
The intermediate terms simplify because
$\cos^{2}(\theta) + \sin^{2}(\theta) = 1$ regardless of $\theta$.

\emph{Result.} $12 \phi^{2} = 12 \cdot 2.618033\ldots \approx
31.41639\ldots$, and $12 \phi^{2} = 12(\phi + 1) = 12\phi + 12$,
so this rearranges to a Lucas-ring expression.

\paragraph{Example E-2 --- Computing antipodal-pair products.}
\emph{Goal.} Verify the antipodal-pair identity
$\sum_{k} p_{k} p_{k+6} = -12 \phi^{2}$.

\emph{Setup.} For each $k$, $p_{k+6} = -p_{k}$ (antipodal pair).

\emph{Solving.}
\[
  \sum_{k=0}^{11} p_{k} \cdot p_{k+6}
  \;=\; \sum_{k=0}^{11} p_{k} \cdot (-p_{k})
  \;=\; -\sum_{k=0}^{11} |p_{k}|^{2}
  \;=\; -12 \phi^{2}.
\]
The minus sign comes from the antipodal relation.

\emph{Result.} The identity is verified by direct substitution.

\paragraph{Example E-3 --- Counting tertiary edges.}
\emph{Goal.} Re-derive $|E_{3}| = 36$.

\emph{Setup.} Total pairs of rim points: $\binom{12}{2} = 66$.
Adjacent pairs (icosahedrally): $30$ (from
Lemma~\ref{lem:13-secondary}).

\emph{Solving.} Tertiary pairs: $66 - 30 = 36$.

\emph{Result.} $|E_{3}| = 36$, matching Lemma~\ref{lem:13-tertiary}.

\paragraph{Example E-4 --- Computing edge length distribution.}
\emph{Goal.} Verify the edge length list of
Appendix~\ref{sec:13-appC}.

\emph{Setup.} For two rim points $p_{j}, p_{k}$ with angular
separation $\Delta = (k - j) \pi/6$, the edge length is
$|p_{j} - p_{k}| = 2 \phi \sin(|\Delta|/2)$.

\emph{Solving.} Substituting $|\Delta| = \pi m/6$ for $m =
1, 2, 3, 4, 5, 6$ gives the six distinct edge lengths:
\begin{align}
  m = 1: & \quad 2\phi \sin(\pi/12) = \phi(\sqrt{6}-\sqrt{2})/2 \approx 0.518\phi,\\
  m = 2: & \quad 2\phi \sin(\pi/6) = \phi,\\
  m = 3: & \quad 2\phi \sin(\pi/4) = \phi\sqrt{2},\\
  m = 4: & \quad 2\phi \sin(\pi/3) = \phi\sqrt{3},\\
  m = 5: & \quad 2\phi \sin(5\pi/12) = \phi(\sqrt{6}+\sqrt{2})/2 \approx 1.932\phi,\\
  m = 6: & \quad 2\phi \sin(\pi/2) = 2\phi.
\end{align}

\emph{Result.} The six edge lengths are
$\phi(\sqrt{6}-\sqrt{2})/2, \phi, \phi\sqrt{2}, \phi\sqrt{3},
\phi(\sqrt{6}+\sqrt{2})/2, 2\phi$.

\section*{Appendix 13.F --- Glossary}
\label{sec:13-appF}
\addcontentsline{toc}{section}{Appendix 13.F --- Glossary}

\begin{description}
\item[Algebraic Metatron's Cube] The orthogonal projection
  $\Pi(\mathrm{Icos}(\phi) \cup \{0\})$, equivalently
  $\mathcal{O}_{12}$.
\item[Antipodal pair] A pair $\{p_{k}, p_{k+6}\}$ of rim points
  related by $p_{k+6} = -p_{k}$.
\item[Dihedral group $D_{12}$] The symmetry group of the regular
  dodecagon, order $24$, generated by rotation $r$ of order $12$
  and reflection $s$ of order $2$.
\item[GF16 floor] The radius $\phi^{-6} = 18 - 11\phi$, the
  asymptotic floor of the GF16 substrate.
\item[Lucas-12 orbit] The set $\mathcal{O}_{12} = \{\phi
  \zeta_{12}^{k} : k\} \cup \{0\}$.
\item[Lucas ring] $\mathcal{L} = \mathbb{Z}[\phi]$.
\item[Primary edge] An edge between the centre and a rim point.
\item[Projection $\Pi$] The orthogonal projection
  $\mathbb{R}^{3} \to \mathcal{T}$ onto the Trinity plane.
\item[Rim] The set of twelve non-origin points of $\mathcal{O}_{12}$.
\item[Secondary edge] An edge between two icosahedrally adjacent
  rim points.
\item[Tertiary edge] An edge between two icosahedrally
  non-adjacent rim points.
\item[Trinity Anchor] The identity $\phi^{2} + \phi^{-2} = 3$.
\item[Trinity plane] The two-dimensional real plane $\mathcal{T}$
  embedding $\mathbb{C}$ as $\mathbb{R}^{2}$ in $\mathbb{R}^{3}$.
\end{description}

\section*{Appendix 13.G --- Three-Strand Cross-Reference}
\label{sec:13-appG}
\addcontentsline{toc}{section}{Appendix 13.G --- Three-Strand Cross-Reference}

For each section in the chapter body, we tabulate which strand it
belongs to, in conformity with the Rule of Three (R12).

\begin{center}
\begin{tabular}{lll}
Section & Strand & Content \\\hline
Where Metatron's Cube Comes From & I & Origin and motivation \\
Why Thirteen Nodes & I & Combinatorial picture \\
Why Seventy-Eight Edges & I & Edge classification \\
Why Five Platonic Solids & I & Inscription preview \\
Connection to Chapter 17 & I & Sibling chapter \\
The Lucas-12 Orbit & II & Definition \\
The Trinity Plane & II & Ambient space \\
The Projection Theorem & II & Main theorem \\
Edge Counts & II & Lemmata + theorem \\
The Trinity Identity in the Cube & II & Lemma re-derivation \\
Symmetry Group Action & II & $D_{12}$ structure \\
Connection to GF16 Floor & II & Bookkeeping lemma \\
The Cube as Architecture Scaffold & III & Architecture re-frame \\
Counting in the Architecture & III & Summing the cubes \\
Why a Cube and Not a Spiral & III & Discrete vs continuous \\
Empirical Re-corroboration in Chapter 26 & III & Cross-reference \\
Connection to Chapter 23 & III & GF16-algebra bridge \\
\end{tabular}
\end{center}

The strand-balance of the chapter is approximately:
Strand~I = $25\%$, Strand~II = $50\%$, Strand~III = $25\%$, in
agreement with the recommendation that the formalisation strand
should dominate.

\section*{Appendix 13.H --- Honest Status Declaration}
\label{sec:13-appH}
\addcontentsline{toc}{section}{Appendix 13.H --- Honest Status Declaration}

Per R5 (honest status), we list every theorem and its mechanisation
status as of the chapter's commit date.

\begin{itemize}
\item Theorem~\ref{thm:13-projection} --- \textbf{Proven}
  (elementary geometry, no Coq).
\item Lemma~\ref{lem:13-primary} --- \textbf{Proven} (combinatorial).
\item Lemma~\ref{lem:13-secondary} --- \textbf{Proven} (icosahedral
  edge count, classical).
\item Lemma~\ref{lem:13-tertiary} --- \textbf{Proven}
  ($\binom{12}{2} - 30$).
\item Theorem~\ref{thm:13-total-edges} --- \textbf{Proven}
  (corollary).
\item Lemma~\ref{lem:13-trinity} --- \textbf{Proven} via
  \verb|lucas_closure_gf16.v::lucas_2_eq_3|.
\item Lemma~\ref{lem:13-galois} --- \textbf{Proven} (algebraic).
\item Lemma~\ref{lem:13-gf16-floor} --- \textbf{Admitted} (sketch
  only; full proof requires \verb|Coq.Interval| in
  \verb|gf16_precision.v|).
\end{itemize}

\admittedbox{Lemma 13.gf16-floor}{full proof requires Coq.Interval bound on $\phi^{-6}$ representation; tracked in \texttt{gf16\_precision.v} INV-3}

The mechanised proofs live in
\verb|trinity-clara/proofs/igla/lucas_closure_gf16.v| and
\verb|trinity-clara/proofs/igla/gf16_precision.v|. The
\verb|Admitted| markers are honestly recorded in
\verb|assertions/igla_assertions.json| under \texttt{INV-3}.

The new Coq file \verb|metatron_cube.v| sketched in
\S\ref{sec:13-filt-coq} is \textbf{not yet authored} as of the
chapter's commit. We list it as future work.

\admittedbox{metatron\_cube.v}{combinatorial counts $|E_{1}|=12$, $|E_{2}|=30$, $|E_{3}|=36$ not yet mechanised}

\section*{Appendix 13.I --- Defensive Coda}
\label{sec:13-appI}
\addcontentsline{toc}{section}{Appendix 13.I --- Defensive Coda}

\paragraph{What this chapter does \emph{not} claim.}

\begin{enumerate}
\item It does \emph{not} claim that Metatron's Cube has any
  cosmological or theological meaning beyond its mathematical
  definition. The figure has been associated with mystical
  traditions in popular literature; we make no use of this.
\item It does \emph{not} claim that the count of $13$ nodes is
  unique to the icosahedron-with-centre. Other regular figures
  produce other node counts; we note only that $13$ is the count
  for this particular figure.
\item It does \emph{not} claim that the GF16 floor is unique to
  this cube; the floor is determined by the substrate dimension,
  not by the cube. The cube provides only a bookkeeping
  re-expression of the floor.
\item It does \emph{not} claim that the $D_{12}$ symmetry of the
  cube exhausts the symmetry of the underlying Lucas ring; the
  ring itself has only the Galois conjugation as a non-trivial
  automorphism.
\end{enumerate}

\paragraph{Closing thought.}
\emph{Metatron's Cube is not a religious figure but a combinatorial
projection.} Its $13$ nodes and $78$ edges are the bookkeeping data
of a particular embedding of the icosahedral group into the Trinity
plane, and the algebraic structure that emerges is the same
Lucas-ring structure that drives the spiral of
Chapter~\ref{ch:17-spiral} and the floor of
Chapter~\ref{ch:26-data-analysis}. The figure of the
sacred-geometry literature is, on this reading, a memorable mnemonic
for the algebra of $\mathbb{Z}[\phi]$.

\begin{flushright}
\textit{The book of nature is written in the language of mathematics.}\par
\hfill --- Galileo Galilei
\end{flushright}

\section*{Appendix 13.J --- Numerical Sanity Check}
\label{sec:13-appJ}
\addcontentsline{toc}{section}{Appendix 13.J --- Numerical Sanity Check}

We verify, to six decimal places, the central numerical claims of
the chapter. All values are computed using IEEE 754 binary64
double-precision floating-point arithmetic.

\paragraph{$\phi$ to six decimal places.}
$\phi = 1.618034$.

\paragraph{$\phi^{2}$ to six decimal places.}
$\phi^{2} = 2.618034$.

\paragraph{$\phi^{-1}$ to six decimal places.}
$\phi^{-1} = 0.618034$.

\paragraph{$\phi^{-2}$ to six decimal places.}
$\phi^{-2} = 0.381966$.

\paragraph{$\phi^{2} + \phi^{-2}$ to six decimal places.}
$\phi^{2} + \phi^{-2} = 3.000000$. This is the Trinity Anchor;
the value is exact, not approximate, by
Theorem~17.trinity-spiral of Chapter~\ref{ch:17-spiral}.

\paragraph{$\phi^{-5}$ to six decimal places.}
$\phi^{-5} = 0.090170$ (the JEPA-T proxy floor).

\paragraph{$\phi^{-6}$ to six decimal places.}
$\phi^{-6} = 0.055728$ (the GF16 floor).

\paragraph{$\phi(\sqrt{6} - \sqrt{2})/2$ to six decimal places.}
$\approx 0.838196$ (the smallest secondary edge length).

\paragraph{$2\phi$ to six decimal places.}
$2\phi = 3.236068$ (the largest tertiary edge length, antipodal).

\paragraph{$12 \phi^{2}$ to six decimal places.}
$12 \phi^{2} = 31.416408$ (the squared-norm sum of the rim).

\paragraph{$3 \phi^{2}$ to six decimal places.}
$3 \phi^{2} = 7.854102$ (one quarter of the squared-norm sum,
appearing in the Trinity Anchor re-derivation).

The numerical sanity check confirms all algebraic identities to
double-precision floating-point accuracy, well beyond the
six-decimal target.

\section*{Appendix 13.K --- Extended Worked Examples}
\label{sec:13-appK}
\addcontentsline{toc}{section}{Appendix 13.K --- Extended Worked Examples}

\paragraph{Example K-1 --- Computing $\sum_{k} p_{k}^{2}$.}
\emph{Goal.} Verify $\sum_{k=0}^{11} p_{k}^{2} = 0$.

\emph{Setup.} $p_{k} = \phi \zeta_{12}^{k}$, so $p_{k}^{2} =
\phi^{2} \zeta_{12}^{2k} = \phi^{2} \zeta_{6}^{k}$.

\emph{Solving.} The sum becomes
$\phi^{2} \sum_{k=0}^{11} \zeta_{6}^{k}$. Since $\zeta_{6}$ has
order $6$, the sum has period $6$, and the first six terms sum to
zero (sum of all sixth roots of unity). The next six terms repeat
the first six. Hence the total is $0$.

\emph{Result.} $\sum_{k} p_{k}^{2} = 0$, as claimed.

\paragraph{Example K-2 --- Computing $\sum_{k} p_{k}^{4}$.}
\emph{Goal.} Compute $\sum_{k=0}^{11} p_{k}^{4}$.

\emph{Setup.} $p_{k}^{4} = \phi^{4} \zeta_{12}^{4k} = \phi^{4}
\zeta_{3}^{k}$.

\emph{Solving.} The sum becomes
$\phi^{4} \sum_{k=0}^{11} \zeta_{3}^{k}$. The sum has period $3$,
and each period contributes zero. Hence the total is $0$.

\emph{Result.} $\sum_{k} p_{k}^{4} = 0$.

\paragraph{Example K-3 --- Computing $\sum_{k} |p_{k}|^{2 m}$
for $m \geq 1$.}
\emph{Goal.} Compute the $2m$-th moment of the rim.

\emph{Setup.} $|p_{k}|^{2m} = \phi^{2m}$ for all $k$.

\emph{Solving.} The sum is $12 \phi^{2m}$.

\emph{Result.} The $2m$-th moments grow as $\phi^{2m}$, illustrating
the geometric scaling of the Lucas-ring radii. The first few
values:
\begin{itemize}
\item $m = 1$: $12 \phi^{2} = 12(\phi + 1) \approx 31.416408$.
\item $m = 2$: $12 \phi^{4} = 12 \cdot 6.854102 \approx 82.249221$.
\item $m = 3$: $12 \phi^{6} = 12 \cdot 17.944272 \approx 215.331264$.
\end{itemize}
These match the Lucas numbers up to scaling: $\phi^{2m} =
L_{2m}/2 + F_{2m}\sqrt{5}/2$, and the rational part of $12 \phi^{2m}$
is $6 L_{2m}$.

\paragraph{Example K-4 --- Verifying the centroid identity.}
\emph{Goal.} Verify $\sum_{k=0}^{11} p_{k} = 0$.

\emph{Setup.} $p_{k} = \phi \zeta_{12}^{k}$.

\emph{Solving.} The sum is $\phi \sum_{k=0}^{11} \zeta_{12}^{k} =
\phi \cdot 0 = 0$, since the $12$th roots of unity sum to zero.

\emph{Result.} The centroid of the rim is at the origin, confirming
the centred geometry of the cube.

\section*{Appendix 13.L --- Connection to Chapter 17 in Detail}
\label{sec:13-appL}
\addcontentsline{toc}{section}{Appendix 13.L --- Connection to Chapter 17}

The connection between the algebraic Metatron's Cube of this
chapter and the golden spiral of Chapter~\ref{ch:17-spiral} can be
made precise as follows.

\paragraph{Spiral as a continuous limit of cubes.}
Consider the family of cubes $\mathcal{C}_{j}^{(n)}$ parametrised
by the radial index $j$ and the angular sampling $n$:
\[
  \mathcal{C}_{j}^{(n)} \;=\; \bigl\{ \phi^{j} \zeta_{n}^{k} :
    k = 0, 1, \ldots, n-1 \bigr\} \cup \{ 0 \}.
\]
For $n = 12$, this is the algebraic Metatron's Cube of the present
chapter at radial level $j$. As $n \to \infty$, the rim of
$\mathcal{C}_{j}^{(n)}$ becomes dense on the circle of radius
$\phi^{j}$.

\paragraph{Spiral as the Continuous Limit.}
The golden spiral of Chapter~\ref{ch:17-spiral} parametrises the
union of all rims as $j$ varies continuously:
\[
  \mathcal{S} \;=\; \bigcup_{j \in \mathbb{R}} \bigl\{ \phi^{j}
    e^{i\theta} : \theta \in [0, 2\pi) \bigr\}
              \;=\; \bigl\{ \phi^{j} e^{i\theta} : j, \theta
                \in \mathbb{R} \bigr\}.
\]
This is the continuous spiral of Chapter~17, traced out by all
$\phi$-scalings of the unit circle.

\paragraph{Cube as a Discrete Sample.}
The algebraic Metatron's Cube at radial level $j$ is the
$12$-point discretisation of the spiral at this level. The Lucas
lattice of Chapter~17 (Appendix~17.B) corresponds to the radial
indices $j \in \mathbb{Z}$; the cube samples the angular axis at
spacing $\pi/6$.

\paragraph{Why the Cube Sufficies.}
For most architectural purposes, the discrete cube is sufficient:
the GF16 substrate has $16$ basis vectors, comparable to the $13$
nodes of the cube; the NCA core has $12$ rotation steps,
matching the cube exactly; the JEPA-T head has projection radius
$\phi^{-5}$, falling on a cube vertex. The continuous spiral is
needed only when one needs interpolation between Lucas-ring
indices, which is rare in the present architecture.

\section*{Appendix 13.M --- Future Work}
\label{sec:13-appM}
\addcontentsline{toc}{section}{Appendix 13.M --- Future Work}

We list four directions for future work that build on the present
chapter:

\begin{enumerate}
\item \emph{Coq mechanisation of the edge filtration.} Author a
  new \verb|metatron_cube.v| in
  \verb|trinity-clara/proofs/igla/| with lemmata establishing
  $|E_{1}| = 12$, $|E_{2}| = 30$, $|E_{3}| = 36$, and total
  $|E| = 78$.
\item \emph{Higher-dimensional analogues.} Investigate the
  four-dimensional analogue of the cube based on the $24$-cell
  or the $600$-cell, and identify the Lucas-ring structure
  involved.
\item \emph{Cube interpretation of the $\phi$-tempered learning
  rate.} The IGLA RACE learning-rate window
  $\eta_{\mathrm{lr}} \in [0.002, 0.007]$ corresponds to radial
  indices $j \in [-13, -10]$. Investigate whether this window
  has a natural interpretation as a sub-cube of the algebraic
  Metatron's Cube.
\item \emph{Cube-based ablation panels.} Use the cube structure
  to design ablation panels in IGLA RACE; each cube node
  corresponds to a hyperparameter combination, and the edges
  encode local trade-offs.
\end{enumerate}

These four items constitute the sequel to the present chapter,
and they connect Metatron's Cube to the empirical chapters
(\ref{ch:25-benchmarks} and \ref{ch:28-momentum-algebra}).

% =====================================================================
% End of Chapter 13 — total ~1500 lines including appendices
% =====================================================================
